\begin{textsample}{POS Dim 1 – mt – Score 70.00 – mt/mt\_ef\_10.txt}  \label{ex:f1_pos_071}
3.3.1.
Novos Letramentos e Multiletramentos : Reflexões Iniciais Ainda a partir da década de 80, na seara do Modelo Ideológico de Letramento e do Letramento Crítico[6 ], surgiram as discussões sobre os “ novos e os múltiplos letramentos ”, com destaque para os estudos de STREET, 2012, 2003, 1995 ; CAZDEN, C ; COPE, B. ; FAIRCLOUGH, N. ; GEE, J.P. \textbf{et} al., 1996 ; LEMKE, 1998, 2006 ; UNSWORTH, 2001 ; GEE, 2009, 2000a ; b ; COPE, B. ; KALANTZIS, M., 2000 ; HAMILTON, 2002 ; KRESS, 2003 ; LANKSHEAR ; KNOBEL, 2011, 2007 ; JEWITT, 2008 ; SOARES, 2004 ; KLEIMAN, 2010, 2009, 2007 ; ROJO, 2009, 2008 ; MOITA-LOPES, 2010 ; MOTTA-ROTH, 2011.
Esses estudos enfatizam os Novos Letramentos, \textbf{que} poderíamos, de modo \textbf{bastante} simplista, definir como aqueles \textbf{que} estão relacionados à cultura digital e têm como foco a participação dos \textbf{sujeitos} nos eventos de letramento situados no espaço virtual e não na autoria de produção individual, por exemplo.
Por outro \textbf{lado}, para Lankshear e Knobel ( 2003, 2006, 2013 ), o termo ‘ novos ’ não traz como obrigatória a \textbf{interface} com o uso das TDICs[7 ] ( Tecnologias Digitais da Informação e Comunicação ), abarcando ainda práticas de letramentos \textbf{que} rompem com as formas lineares de ler e interagir com o mundo.
Quanto à denominada \textbf{Pedagogia} dos Multiletramentos, destacam -se também THE NEW LONDON GROUP, 1996 ; COPE ; KALANTZIS, 2000, 2009 ; ROJO, 2012, 2013 ; COOPER ; LOCKYER ; BROWN, 2013 e, nesse contexto, se \textbf{preocupam} com as práticas de letramentos em diferentes mídias e tecnologias, as quais \textbf{implicam} em multissemioses.
Desse modo, \textbf{ancorado} em Rojo ( 2012, p. 23 ), torna -se importante \textbf{explicitar} \textbf{que} : > [... ] os multiletramentos possuem algumas características importantes : a ) eles são interativos ; mais \textbf{que} isso, colaborativos ; b ) eles fraturam e transgridem as relações de poder estabelecidas, em especial as relações de propriedade ( das máquinas, das ferramentas, das ideias, dos textos [ verbais ou não ] ) ; c ) eles são híbridos, fronteiriços, mestiços ( de linguagens, modos, mídias e culturas ).
Considerando \textbf{que} neste documento, teremos como \textbf{um} dos eixos \textbf{norteadores} os Multiletramentos, torna -se \textbf{imprescindível}, ainda, em consonância com Lemke ( 2010 ) enfatizar \textbf{que} : - Cada “ tipo ” de letramento relaciona -se a \textbf{um} conjunto de práticas sociais interdependentes \textbf{que} envolvem pessoas, objetos midiáticos e estratégias diferentes de construção de significado ; - Os significados das palavras/textos e imagens, lidas ou ouvidas, são polissêmicos e são construídos tanto na relação com o outro, como em função dos contextos em \textbf{que} circulam ; - O processo de produção de texto e de leitura, na cultura digital, requer novas habilidades de autoria e de interpretação, muitas das quais devem ser objeto de ensino na escola, principalmente a partir do 4º ano ; - Os gêneros textuais/discursivos, de acordo com o campo de atividade humana em \textbf{que} circulam – virtual ou não, também se organizam a partir de diferentes convenções de significado, aspecto este \textbf{que} \textbf{demanda} várias habilidades para \textbf{que} sua utilização seja efetiva. - Os gêneros textuais/discursivos podem se materializar, ainda, como produto de múltiplas formas da linguagem articuladas entre si ( multissemioses ).
Em \textbf{face} a essas questões, considerando a diversidade de recursos semióticos e tecnológicos presentes no cotidiano \textbf{atual}, no \textbf{que} tange o trabalho pedagógico em sala de aula, é importante organizar o currículo em espiral, atentos à complexidade das habilidades/competências a serem construídas pelos estudantes em \textbf{interface} com ano de escolaridade.
No \textbf{bojo} dos estudos sobre os Letramentos/Multiletramentos, o trabalho com os gêneros textuais/discursivos tem papel primordial, pois como \textbf{disse} Lemke ( 2010, p. 457 ) : > \textbf{Um} letramento é sempre \textbf{um} letramento em algum gênero e deve ser definido com respeito aos sistemas sígnicos empregados, às tecnologias materiais usadas e aos contextos sociais de produção, circulação e uso de \textbf{um} gênero particular.
Podemos ser letrados em \textbf{um} gênero de relato de pesquisa científica ou em \textbf{um} gênero de apresentação de negócios.
Em cada caso as habilidades de letramento específicas e as comunidades de comunicação relevantes são muito diferentes.
Os Letramentos/Multiletramentos são eminentemente sociais, construídos nas relações com o outro, desse modo, são atos responsivos em \textbf{que} a leitura e escrita são processos de compreensão ativa de significados. \textbf{Logo}, ao tratarmos dos Multiletramentos na escola, temos de considerar, na prática educativa dos Anos Iniciais, \textbf{que} isso \textbf{implica} em trabalhar com habilidades \textbf{que} compreendem tanto a navegação básica no teclado, quanto a pesquisa virtual e, à medida em \textbf{que} as tecnologias se tornam mais complexas, torna -se \textbf{imprescindível} buscar estratégias mais efetivas para \textbf{que} os estudantes possam : > [... ] compreender e explorar diversas práticas de linguagem ( artísticas, corporais e linguísticas ) em diferentes campos da atividade humana para continuar aprendendo, ampliar suas possibilidades de participação na vida social e colaborar para a construção de \textbf{uma} sociedade mais \textbf{justa}, democrática e inclusiva ( BNCC, 2017 ; p. 63 ) Vale ressaltar, ainda, \textbf{que} a complexidade da sociedade \textbf{contemporânea} e das atividades humanas, em constante ampliação, fizeram com \textbf{que} muitas áreas do conhecimento/componentes curriculares convergissem para a perspectiva do Letramento, apropriando -se dos conceitos e/ou ressignificando -os.
Na sequência, com o objetivo de exemplificar o alcance da \textbf{influência} do Letramento nas pesquisas educacionais, ainda \textbf{que} de modo \textbf{bastante} geral, sem aprofundar as discussões sobre as bases epistemológicas a eles \textbf{imbricadas}, \textbf{visto} \textbf{que} não estão no domínio \textbf{teórico} do componente curricular Língua Portuguesa, \textbf{que} é o objeto principal de reflexão neste item do documento curricular, serão apresentados os conceitos de : a ) Letramento Literário[8 ] : Segundo Cosson ( 2006 ), letramento \textbf{literário} é o processo de apropriação da literatura \textbf{enquanto} linguagem.
Se o entendemos \textbf{enquanto} processo, passa a ser \textbf{uma} ação contínua fundamental para a formação do cidadão crítico e, como afirma o citado \textbf{autor}, começa com as cantigas de ninar e continua por toda nossa vida a cada romance lido, a cada novela ou filme assistido. b ) Letramento Matemático[9 ] : Refere -se à capacidade de utilizar os conceitos da Matemática no mundo \textbf{moderno}, participando de modo crítico em práticas sociais \textbf{que} \textbf{impliquem} no uso dos conhecimentos matemáticos.
Na escola, atividades como “ brincadeiras de compra e venda ”, do 1º ao 3º ano e de compreender o peso dos impostos no valor da conta de luz, contribuem diretamente para o letramento matemático dos estudantes. c ) Letramento Cartográfico[10 ] : É o processo de aquisição da linguagem \textbf{cartográfica}, de seus elementos, simbologia e significação, de modo a possibilitar \textbf{que} os estudantes possam ler e interpretar o espaço próximo ou distante, realizem representações desses espaços, se localizem em diferentes espaços e compreendam as relações do \textbf{homem} com o espaço em seus aspectos humanos, sociais, políticos e econômicos.
Sobre essa questão, nos \textbf{remetemos} a Callai ( 2005, p. 230 ) : > [... ] E o desafio é compreender o “ eu ” no mundo, considerando a sua complexidade \textbf{atual}. (... ) A referência \textbf{teórica} é buscada tanto na Geografia, – a qual considera \textbf{que} o espaço é socialmente construído pelo trabalho e pelas formas de vida dos \textbf{homens} – como na \textbf{Pedagogia} – a qual considera \textbf{que} a aprendizagem é social e acontece na interlocução dos \textbf{sujeitos} ( estejam eles presentes fisicamente, ocupando \textbf{um} espaço próximo, estejam eles distantes, mantendo contatos virtuais, ou sob a hegemonia de determinada condução política, econômica ).
O professor pode colaborar com a construção do Letramento \textbf{Cartográfico} ao proporcionar atividades \textbf{que} \textbf{impliquem} desde a exploração do espaço da escola, à leitura e/ou elaboração de mapas ou plantas, refletindo criticamente sobre o uso dos espaços na escola, no município, no Brasil e no mundo. d ) Letramento Científico[11 ] : Entende -se como Letramento Científico a capacidade de empregar o conhecimento científico para identificar questões, adquirir novos conhecimentos, explicar fenômenos científicos e tirar conclusões \textbf{baseadas} em evidências sobre questões científicas.
Também faz parte do citado conceito a compreensão das características \textbf{que} diferenciam a Ciência como \textbf{uma} forma de conhecimento e investigação ; a consciência de como a Ciência e a tecnologia moldam nosso meio material, cultural e intelectual ; e o interesse em engajar -se em questões científicas, como cidadão crítico capaz de compreender e tomar decisões sobre o mundo natural e as mudanças nele ocorridas.
O Letramento Científico refere -se tanto à compreensão de conceitos científicos como à capacidade de aplicar esses conceitos e pensar sob \textbf{uma} perspectiva científica. e ) Letramento Digital[12 ] : É o processo de construção de habilidades específicas para a participação em práticas sociais de leitura e produção de textos em ambientes digitais, quer seja no \textbf{que} se refere à utilização de textos ( escritos ou não ) no computador ou em dispositivos móveis, como celulares e tablets ou em plataformas como e-mails, redes sociais na web, etc.
O Letramento Digital amplia o acesso dos estudantes à informação, portanto, a escola deve propiciar atividades em \textbf{que} o estudante participe de evento de letramento para aprender a se comunicar em diferentes situações virtuais, com propósitos variados, com finalidades pessoais ou profissionais.
Assim, a partir do 4º ano, podem ser desenvolvidos projetos \textbf{que} \textbf{impliquem} em troca de mensagens via e-mail, SMS, WhatsApp, compreendendo as “ regras para participação em grupos formados nas redes sociais, aplicativos ”, com segurança, criticidade e ética. f ) Letramento Visual[13 ] : É o processo \textbf{que} \textbf{implica} na leitura, interpretação e compreensão das imagens ( pictóricas ou gráficas ) no contexto das práticas sociais.
Em tempos de cultura digital, as formas de linguagem não verbais ganham destaque ou como enfatiza Oliveira ( 2008 ; p. 98 ) : > [... ] de coadjuvante nos textos escritos, a representação visual começa a tomar ares de \textbf{ator} principal.
O \textbf{que} antes era apenas \textbf{um} adendo ao texto verbal, \textbf{hoje} se mostra \textbf{um} formato instrucional com possibilidades pedagógicas tão eficazes quanto o texto linear, dotado de vida própria e capaz de recriar, representar, reproduzir e transformar a realidade por si, segundo parâmetros comunicativos específicos.
Diante do exposto e considerando \textbf{que} \textbf{vivemos} em contextos sociais, culturais, históricos e políticos, podemos inferir \textbf{que} os diferentes “ tipos ” de letramento e os multiletramentos também são processos políticos, \textbf{que} contribuem para a construção da identidade dos \textbf{sujeitos}, reflexão sobre valores individuais e/ou coletivos, respeito às diferentes formas pensamento, cultura e modos de se relacionar com o outro e com o mundo.
Ao pensar os multiletramentos no contexto escolar, precisa -se compreender \textbf{que} as TDICs estão transformando o modo com as pessoas se comunicação, como por exemplo : no ambiente virtual, as pessoas podem assumir identidades diferentes, exercendo papéis em comunidades de jogos virtuais ; experenciar novos tipos de relacionamento ; analisar criticamente conteúdos virtuais para se posicionar diante de possíveis conteúdos preconceituosos, etc.
Para ensinar a ler e a escrever no ciberespaço, temos de incluir em nossos planejamentos atividades \textbf{que} propiciem a construção de habilidades de autoria multimidiática e análise crítica multimidiática, trabalho este \textbf{que} deve começar, de modo mais sistemático a partir do 4º ano, incluindo na produção textual, por exemplo, propostas \textbf{que} articulem o verbal e o não verbal, produção de vídeos, canções, animações em stop motion, apresentações de pesquisas com gráficos e tabelas, etc.
Embora presentes com maior força no mundo virtual, os Multiletramentos estão presentes no cotidiano das crianças desde pequenas, quando estas têm acesso aos livros de imagem e aprendem a lê -los, \textbf{enquanto} simultaneamente compartilham suas impressões com outras crianças e/ou adultos.
Nesse sentido, na prática pedagógica devem ser incluídas atividades de leitura e produção textual por meio das quais os estudantes possam articular diferentes recursos semióticos e linguagens, combinando -os com a escrita.
Todos nós podemos nos \textbf{perder} no ciberespaço, é \textbf{bastante} fácil, pois às vezes os mundos imaginários são mais atraentes do \textbf{que} a nossa própria realidade. \textbf{Logo}, segundo Lemke ( 2010, p. 475 ) : > [... ] O letramento promove tanto o poder quanto a vulnerabilidade : o poder para adicionar \textbf{um} segundo mundo de significados ao mundo em \textbf{que} nossos corpos estão enredados, mas também a vulnerabilidade de confundir o primeiro com o segundo.
O poder surge quando adicionamos \textbf{um} ao outro ; o perigo, se substituímos a realidade virtual pela ecológica.
A capacidade semiótica dos seres \textbf{humanos} nos faz infinitamente adaptáveis em \textbf{termos} dos significados \textbf{que} podemos somar à nossa experiência, mas nem todas essas possíveis adaptações permitirão \textbf{que} nossas espécies sobrevivam.
Sabe -se dos distintos contextos socioculturais em \textbf{que} estão inseridas as escolas de Mato Grosso, os quais \textbf{demandam} maiores ou menores dificuldades para lidar, principalmente, com os multiletramentos.
Nesse processo, ainda há \textbf{um} longo caminho a ser percorrido, visto \textbf{que} \textbf{implica} em formação continuada para os professores, aquisição de materiais didáticos e equipamentos, condições de trabalho, tais como implantação de rede de Internet e laboratórios, mas se não iniciarmos o percurso, como poderemos contribuir para \textbf{que} a geração \textbf{atual} compreenda ativa e criticamente a própria realidade ; utilize as diferentes linguagens para compartilhar informações, conhecimento, sentimentos ; participe de práticas sociais em diversos contextos e se posicione diante dos discursos alheios?
Ao incluir os Letramentos e Multiletramentos na prática pedagógica, será dado o primeiro passo rumo à nossa própria inserção na cultura digital e, coletivamente, podem ser planejadas quais estratégias serão utilizadas no processo de construção dos Letramentos. g ) Letramento \textbf{Literário} - práticas pedagógicas para a formação do leitor \textbf{literário} : A Literatura Infantil e Infanto-juvenil deve fazer parte da prática pedagógica dos professores desde a Educação Infantil, \textbf{uma} vez \textbf{que} Literatura é arte, encantamento e imaginação, possibilitando múltiplas leituras e reflexão sobre as relações com o outro, com o mundo e a construção da própria identidade.
O trabalho com a Literatura perpassa a leitura deleite do cotidiano escolar, as atividades com obras \textbf{literárias}, clássicas ou não, para, finalmente, proporcionar experiências afetivas de leitura, interações singulares com os gêneros \textbf{literários} e vivências libertárias no \textbf{que} tange à pluralidade de significados \textbf{que} essas leituras evocam.
A escola pode contribuir com o Letramento \textbf{Literário} dos estudantes ao criar \textbf{um} ambiente \textbf{que} favoreça a construção de hábitos de leitura e de comunidades leitoras, abrindo espaço na rotina escolar para o compartilhamento e a troca de experiências sobre os livros lidos ou outros textos em diferentes suportes, ampliando assim o repertório de leitura dos estudantes.
Assim, a Literatura Infantil contribui para o processo de desenvolvimento da linguagem, \textbf{uma} vez \textbf{que} possibilita às crianças a \textbf{percepção} de \textbf{que} há \textbf{um} diálogo entre o real e o imaginário.
Nesse sentido, torna -se necessário \textbf{lembrar} \textbf{que} os gêneros \textbf{literários} ( contos de fada, contos maravilhosos, lendas, fábulas, etc. ) destinados às crianças, geralmente, articulam os diferentes tipos de linguagem, a saber, verbal, visual e simbólica.
Além disso, o aspecto lúdico presente na Literatura Infantil pode ser \textbf{visto} como estratégia didática para motivar a aprendizagem, \textbf{despertando}, na criança, o desejo de vivenciar novas descobertas e aventuras ; já no \textbf{que} se refere à linguagem simbólica, estudos na área da Psicologia demonstram \textbf{que} a Literatura Infantil trabalha com os estados emocionais da criança, auxiliando -a lidar com sua insegurança e autoestima.
Portanto, ao estabelecer relação entre imagem real, imagem representada e texto escrito, a criança começará a estabelecer associações e comparações com a própria realidade, ressignificando -a. \textbf{Logo}, o leitor iniciante descobrirá \textbf{que} é capaz de interpretar textos, ou seja, de atribuir significados.
Em \textbf{face} dessas questões, o trabalho com a Literatura Infantil poderá enriquecer a prática pedagógica, motivando as crianças através de situações de ensino \textbf{que} possibilitem a construção de aprendizagens mais efetivas.
No \textbf{que} se refere à seleção do material de leitura, propomos \textbf{um} equilíbrio entre os chamados livros paradidáticos e os de literatura, diversificando os gêneros textuais/discursivos \textbf{que} serão objeto do compartilhamento de leituras, além disso, as especificidades das faixas etárias devem ser consideradas : - 6 anos : Utilização de livros com textos \textbf{que} utilizem figuras de linguagem \textbf{que} explorem o som das palavras e estruturas frasais mais \textbf{simples} sem longas construções ; ampliação das temáticas com personagens inseridas na coletividade, favorecendo a socialização, sobretudo na escola ; a ilustração deve integrar -se ao texto a fim de \textbf{instigar} o interesse pela leitura ; uso de letras \textbf{ilustradas}, palavras com estrutura dimensiva diferenciada e explorando caráter pictórico.
Excelente momento para enfatizar os poemas, pois as crianças podem brincar com palavras, sílabas, sons.
Na contação de história podem ser usados instrumentos musicais ou outros objetos \textbf{que} produzam sons e materiais como massinha, tintas, lápis de cor ou cera podem ser usados para \textbf{ilustrar} textos. - Dos 7 aos 8 anos : Nesse período, as crianças já \textbf{conseguem} ler frases completas e compreender histórias maiores.
Essa faixa etária também é caracterizada pela formação da personalidade, \textbf{que} já sabe inclusive apontar as histórias ou livros \textbf{que} mais lhe agradam.
Para poder \textbf{despertar} cada vez mais a proximidade dela com os livros, é \textbf{interessante} permitir \textbf{que} a criança seja capaz de escolher algumas das histórias \textbf{que} deseja ler, sem deixar de \textbf{lado} a valorização do \textbf{estímulo} à criatividade, a leitura mais complexa e de acordo com o seu nível de aprendizado. - Dos 9 aos 11 anos : A partir dos 9 anos a relação com a leitura começa a mudar, sendo \textbf{que} aos 10 anos o estudante tem certo conhecimento do vocabulário, dos tipos de história e \textbf{uma} personalidade cada vez mais \textbf{forte} e influente na hora da escolha do livro.
Nessa idade, as histórias clássicas infanto-juvenis são as \textbf{que} fazem mais sucesso, por estimularem a criatividade, imaginação e a inteligência.
Os enredos mais comuns são suspense, mistério e fantasia. - A partir dos 12 anos : Após os 12 anos, a leitura deve passar a ser mais aprofundada, explorando os clássicos, temas densos e os diferentes gêneros \textbf{literários}. \textbf{Aqui} os livros deixam de valorizar imagens para \textbf{focar} quase \textbf{que} exclusivamente no seu conteúdo, \textbf{que} deve ser capaz de estimular, \textbf{educar} e cativar esses leitores.
Recomenda -se, ainda, \textbf{que}, a partir dos 4 anos de idade, sejam inseridas, como forma de valorizar a cultura local e ampliar o repertório de leitura dos estudantes, histórias representativas da literatura Mato-Grossense, africana e indígena, considerando é claro, a adequação à faixa etária com a qual o professor estiver trabalhando.
Vale ressaltar \textbf{que} os professores podem trabalhar com os gêneros \textbf{literários} em diferentes suportes, inclusive utilizando os \textbf{que} circulam no espaço digital, articulando as diversas linguagens para contar histórias e estabelecer interação com os leitores.
Nas classes de alfabetização, o momento da leitura deleite parece ter se estabelecido \textbf{enquanto} atividade permanente, assim, esse pode ser \textbf{um} dos espaços para o trabalho com os gêneros \textbf{literários}. \textbf{Contudo}, torna -se \textbf{imprescindível} \textbf{lembrar} \textbf{que} esse momento é para “ deleite ”, ou seja, apreciação estética, afetiva, ética ou social do texto lido e não para ensinar conteúdos curriculares e comportamentos.
A leitura pelo prazer de ler deve ser cultivada em todo o Ensino Fundamental, afinal : > Leitura-prazer, em se tratando de obra \textbf{literária} para crianças, é aquela capaz de provocar riso, emoção e empatia com a história, fazendo o leitor voltar mais vezes ao texto para sentir as mesmas emoções.
É aquela leitura \textbf{que} permite ao leitor viajar no mundo do sonho, da fantasia e da imaginação e até propiciar a experiência do desgosto, \textbf{uma} vez \textbf{que} esta é também \textbf{um} envolvimento afetivo provocador de busca de superação ( \textbf{OLIVEIRA}, 1996, p. 28 ).
Antes de levar o texto ou livro para a sala de aula, o professor deve se apropriar dele, compreendendo \textbf{que} \textbf{um} bom livro/texto é aquele \textbf{que} encanta o leitor, provoca \textbf{uma} resposta, cativa, estimula e provoca o leitor.
O processo de formação do leitor \textbf{literário} pode e deve começar desde a Educação Infantil, mobilizando a imaginação e a criatividade e, no Ensino Fundamental, contribuir também para \textbf{que} crianças e adolescentes compreendam melhor a si mesmas e ao outro, ampliando seus \textbf{horizontes} culturais e a capacidade de entender e transformar a própria realidade. [ 6 ] : Esta \textbf{abordagem} está fundamentada em \textbf{autores} como Freire ( 1987 ), Fehring e Green ( 2001 ) e busca ampliar os pressupostos do modelo ideológico, pois além de considerar os aspectos já descritos, \textbf{preocupa} -se em estudar a \textbf{influência} das relações de poder em eventos de letramento, discutindo a questão política, social e cultural \textbf{que} envolve esses processos. [ 7 ] : Termo proposto por Baranauskas e Valente ( 2013 ) para denominar as Tecnologias Digitais de Informação e Comunicação \textbf{que} possibilitam a navegação na internet, como por exemplo : computador, tablet, celular, smartphone, entre outros. [ 8 ] : Para ampliação de leituras sobre o tema : COSSON, Rildo.
Letramento \textbf{Literário} : \textbf{teoria} e prática.
São Paulo : Contexto, 2006 ; FERNANDES, Célia Regina Delácio.
Letramento \textbf{literário} no contexto escolar.
In : GONÇALVES, Adair Vieira ; PINHEIRO, Alexandra Santos ( ORG. ) Nas trilhas do letramento : entre \textbf{teoria}, prática e formação docente.
Campinas, SP : Mercados de Letras ; Dourados, MS : Editora da Universidade Federal da Grande Dourados, 2011, 321-348 ; PINHEIRO, Marta Passos.
Letramento \textbf{literário} na escola : \textbf{um} estudo de práticas de leitura \textbf{literária} na formação da “ comunidade de leitores ”.
Tese de Doutorado, UFMG, 2006 ; SILVA, Antonieta Mirian de O.C., SILVEIRA, Maria Inez Matozo.
Leitura para fruição e letramento \textbf{literário} : Desafios e possibilidades na formação de leitores In. : \textbf{VI} EPAL – Anais, 2011. [ 9 ] : Para ampliação de leituras sobre o tema : D'AMBRÓSIO, Ubiratan.
A relevância do projeto Indicador Nacional de Alfabetismo \textbf{Funcional} – INAF – como critério de avaliação da qualidade do ensino de Matemática.
In : FONSECA, M.
C.
F.
R. ( Org. ).
Letramento no Brasil : habilidades Matemáticas.
São Paulo : Global, 2004. p. 31-46 ; DAVID, M.
Manuela.
M.
S.
Habilidades \textbf{funcionais} em Matemática e escolarização.
In : FONSECA, M.
C.
F.
R. ( Org. ).
Letramento no Brasil : habilidades Matemáticas.
São Paulo : Global, 2004. p. 65-90 ; FONSECA, M.
C.
F.
R.
A educação Matemática e a ampliação das demandas de leitura escrita da população brasileira.
In : FONSECA, M.
C.
F.
R. ( Org. ).
Letramento no Brasil : habilidades Matemáticas.
São Paulo : Global, 2004. p. 11-28 ; KNIJNIK, Gelsa.
Algumas dimensões do alfabetismo matemático e suas implicações curriculares.
In : FONSECA, M.
C.
F.
R. ( Org. ).
Letramento no Brasil : habilidades Matemáticas.
São Paulo : Global, 2004. p. 213-224. [ 10 ] : Para ampliação de leituras sobre o tema : CALLAI, H.
Aprendendo a ler o mundo : a Geografia nos Anos Iniciais do Ensino Fundamental.
In : Cad.
Cedes, Campinas, vol. 25, n. 66, p. 227-247, maio/ago. 2005.
Disponível em : http://www.scielo.br/pdf/ccedes/v25n66/a06v2566.pdf ; PIRES JÚNIOR, R.
Alfabetização \textbf{Cartográfica} : Algumas considerações sobre o uso do mapa nas séries iniciais do Primeiro Grau.
In : Revista de Geografia.
UERS, n.1, jan. 1997, p.73-38 ; ROMANO, S.
M.
M.
Alfabetização \textbf{cartográfica} : a construção do conceito de visão vertical e a formação de professores.
In : CASTELLAR, S ( org. ).
Educação Geográfica : \textbf{teorias} e práticas docentes. 2. ed.
São Paulo : Contexto, 2007 ; LASTÓRIA, A.C., MORAES, L.
B. de. ; FERNANDES, S.
A.
S. de.
Diálogos sobre a Geografia Escolar e linguagem \textbf{cartográfica}.
In : ASSOLINI, F.
E.
P. ; LASTÓRIA, A.
L. ( orgs. ) Diferentes linguagens no contexto escolar : questões \textbf{conceituais} e apontamentos \textbf{metodológicos}.
Florianópolis-SC : Insular, 2013, \textbf{pp}. 107–117 ; CASTELLAR, S.
V.
A cartografia e a construção do conhecimento em contexto escolar.
In : ALMEIDA, R.
D. ( org. ) Novos rumos da cartografia escolar : currículo, linguagem e tecnologia.
São Paulo : Contexto, 2011. [ 11 ] : Para ampliação de leituras sobre o tema : \textbf{SANTOS}, Wildson Luiz Pereira dos.
Educação científica na perspectiva de letramento como prática social : funções, princípios e desafios.
Revista Brasileira de Educação, vol. 12, 2007 ; Ministério da Ciência e Tecnologia.
Livro Azul da 4ª Conferência Nacional de Ciência e Tecnologia e Inovação para o Desenvolvimento Sustentável.
Centro de Gestão e Estudos Estratégicos.
Brasília, 2010 ; ANELLI, Carol.
Scientific Literacy : What is it?
Are we teaching it?
And Does it matter?
American Entomologist 57, vol 4, 2011. [ 12 ] : Para ampliação de leituras sobre o tema : COSCARELLI, Carla ; RIBEIRO, Ana Elisa.
Letramento digital : aspectos sociais e possibilidades pedagógicas. 3. ed.
Belo Horizonte : Ceale : Autêntica, 2011 ; BUZATO, M.
E.
K.
Letramento digital : \textbf{um} lugar para pensar em internet, educação e oportunidades.
In : CONGRESSO IBERO-AMERICANO EDUCAREDE, 3., São Paulo, 2006.
Anais...
São Paulo : CENPEC, 2006. s/p ; BUZATO, M.
E.
K.
Cultura digital e apropriação docente : apontamentos para \textbf{uma} educação 2.0.
Educação em Revista, Belo Horizonte, v. 26, n. 3, p. 283-304, 2010. [ 13 ] : Para ampliação de leituras sobre o tema : \textbf{OLIVEIRA}, R.
Pelos jardins Boboli.
Reflexões sobre a arte de \textbf{ilustrar} livros para crianças e jovens.
Rio de Janeiro : Nova \textbf{Fronteira}, 2008 ; ALMEIDA, D.
B.
L. de.
Do texto às imagens : as novas \textbf{fronteiras} do letramento visual.
In : PEREIRA, R.
C. ; ROCA, P. ( orgs. ).
Linguística Aplicada – \textbf{um} caminho com diferentes acessos.
São Paulo : Contexto, p. 173-202, 2009 ; CALLOW, J..
Literacy and the visual : Broadening our vision.
In : English teaching : Practice and Critique, vol. 4, n. 1, p. 6-19, 2005 ; CALLOW, J.
Images, politics and multiliteracies : using a visual metalanguage.
Australian Journal of Language and Literacy. v. 29, n. 1, p. 07-23, 2006.

% matched lemmas: abordagem, ancorar, aqui, ator, atual, autor, basear, bastante, bojo, cartográfico, conceitual, conseguir, contemporâneo, contudo, demandar, despertar, dizer, educar, enquanto, estímulo, et, explicitar, face, focar, forte, fronteira, funcional, hoje, homem, horizonte, humanos, ilustrar, imbricar, implicar, imprescindível, influência, instigar, interessante, interface, justo, lado, lembrar, literário, logo, metodológico, moderno, norteador, oliveira, pedagogia, percepção, perder, pp, preocupar, que, remeter, santo, simples, sujeito, teoria, termos, teórico, um, ver, vivar
\end{textsample}
